\documentclass[journal=jacsat,manuscript=article]{achemso}

\usepackage[version=3]{mhchem}
\usepackage{siunitx}
\usepackage{graphicx}
\usepackage{amsmath}
\usepackage{amssymb}
\usepackage{bm}
\usepackage{booktabs}
\usepackage{array}
\usepackage{hyperref}

\author{Rhys J. Bunting}
\email{bunting4@llnl.gov}
\affiliation[LLNL]
{Quantum Simulations Group, Materials Science Division, Lawrence Livermore National Laboratory, Livermore, California 94550, United States}

\author{Nikhil Rampal}
\affiliation[LLNL]
{Quantum Simulations Group, Materials Science Division, Lawrence Livermore National Laboratory, Livermore, California 94550, United States}

\author{William Engler}
\affiliation[GLOBUS]
{Globus, University of Chicago, Chicago, Illinois 60637, USA}

\author{Ruchika Mahajan}
\affiliation[SLAC]
{SUNCAT Center for Interface Science and Catalysis, SLAC National Accelerator Laboratory, Menlo Park, California 94025, United States}

\author{Kirsten T. Winther}
\affiliation[SLAC]
{SUNCAT Center for Interface Science and Catalysis, SLAC National Accelerator Laboratory, Menlo Park, California 94025, United States}

\author{Michal Bajdich}
\affiliation[SLAC]
{SUNCAT Center for Interface Science and Catalysis, SLAC National Accelerator Laboratory, Menlo Park, California 94025, United States}

\author{Joel B. Varley}
\affiliation[LLNL]
{Quantum Simulations Group, Materials Science Division, Lawrence Livermore National Laboratory, Livermore, California 94550, United States}

\author{Tuan Anh Pham}
\email{pham16@llnl.gov}
\affiliation[LLNL]
{Quantum Simulations Group, Materials Science Division, Lawrence Livermore National Laboratory, Livermore, California 94550, United States}

\title[Online Graph Kinetic Monte Carlo]
{Online Graph Kinetic Monte Carlo}

\abbreviations{KMC, DFT, ASE, NEB}
\keywords{kinetic Monte Carlo, graph theory, surface catalysis, reaction network generation, kinetics}

\begin{document}

\begin{abstract}
Yippeeeee
\end{abstract}

%\section{Introduction}

Kinetic modeling seeks to describe the reactivity and selectivity a catalyst has toward the conversion of reactants to products.
The desired outcome is to theoretically predict these key catalyst properties accurately to either optimize conditions or to find new active and selective catalysts.
There are a variety of methods employed to achieve this.
The simplest method is to use volcano plots, where the adsorption energies of key intermediates are calculated and fitted to a Bell–Evans–Polanyi principle relationship.
This can be combined with linear scaling, giving an optimal adsorption energy for a key intermediate to have the most activity toward a certain product.
With this, no explicit kinetic rate information of elementary reactions on the surface are calculated.
This has been shown to be effective for both heterogeneous catalysis and homogeneous catalysis.
Beyond this method is micro-kinetic modeling, where the kinetic rates of elementary reaction steps on the surface are considered.
A steady state is assumed, where the gradient of the concentration of all surface species (the rate) equals 0.
However, reactions are surface specific and not site specific, causing a reduction in information, with everything considered within a mean-field approximation.
Despite this, surface species concentrations and the surface specific rates of elementary reactions can still be discerned.
For atomic level information of surface reactions, kinetic Monte Carlo offers that resolution, where surface site occupation by species, and subsequent reactions between surface species, are considered.
These elementary steps are discretely modeled as stochastic events.
This allows the reaction to explicitly take place, going beyond just surface concentrations.

The key issue with kinetic modeling is that foreknowledge of the reaction mechanism is required.
An alternative strategy is to consider an exhaustive amount of reaction pathways, ensuring the dominant reaction pathway is safely captured.
The major implication of this is that if new mechanisms are the dominant pathway for new classes of catalysts, the reactivity would be underestimated.
A new method that builds the reaction network lazily, when new reactions on the surface become possible, would radically transform kinetic modeling.
There has been some development to try to achieve this, such as with adaptive KMC, to extend KMC as an on the fly method.
Here, a method is developed that considers possible graph evolutions as reaction classes.
With this method, denoted as online graph kinetic Monte Carlo (OGKMC), reaction products, reactivity, selectivity, and the dominant reaction pathway can be found, with only a catalyst structure, reactants, and reaction conditions as inputs.
%\section{Results and Discussion}

%%%%%%%%
% Fig of covalent radius, structure classification.
%%%%%%%%

OGKMC builds the catalyst and reactant structures as graphs.
The atoms are nodes and bonds are edges.
The nodes store the element, cartesian coordinates, classification, and an occupation flag.
Bonding is classified based on the overlapping elemental covalent radii between two atoms.
In this scenario, edges are built between these nodes that are sufficiently close enough to each other.
As KMC should be performed from states that are local minima, this bonding is classified for structurally minimized structures.
The first step is building the catalyst structure and performing a structural optimization with a potential.
This could be with classical or electronic structure based methods.
After optimization, the bonding is set.
This produces the base graph with nodes and edges.
After this, the reactive part of the catalyst must be designated.
Designating a reactive part of the catalyst enables the complexity of new bonding being formed to be simplified.
This is intuitive: subsurface atoms do not directly bond to reactants; reactions take place almost always at the surface.
In our implementation, any atomic structure can be input, with support for both 2D surfaces and 3D catalysts.
Classification of the surface is through ray casting for 2D surfaces and through generation of a convex hull for 3D catalysts.

%fig.1a -> graph description
%fig.1b -> edge cut off
%fig.1c -> surface classification
%fig.1d -> clique classification
%fig.1e -> graph isomorphs (nshell)
%fig.1f -> graph example of different surface sites.
%fig.1g -> CO adsorbate example

With the surface defined, possible bonding to the surface can be considered.
Possible bonding to the surface is found based on all the geometrically possible coordinations the surface atoms, of varying covalent radii (j + k), have with an adsorbate atom with covalent radius (i).
Each element, with a unique covalent radius, will have different geometrically possible surface coordinations.
For a set of surface atoms $S={1,\ldots,N}$, where surface atom $j$ has position $\mathbf{x}_j$ and covalent radius $r_j$.
This set is from a KD-tree which accelerates the neighbor construction to only viable candidate pairs.
For an adsorbate atom $A$ with covalent radius $r_A$, an edge is placed between surface atoms $j$ and $k$ when $d_{jk}\leq 2r_A+r_j+r_k$, where $d_{jk}$ is the distance between the two surface atoms.
This determines whether a single adsorbate atom can geometrically bond to both surface atoms simultaneously.
With these edges built between all nodes, finding the possible sites is achievable by finding all the cliques within the graph.
A 1-fold clique would be a site that coordinates with a single surface atom (top site); a 2-fold clique is a site that coordinates with two surface atoms (bridge site); and a 3 or 4-fold clique would be a site that coordinates with three or four surface atoms (hollow site).
With efficient modern algorithms for finding cliques in a graph, all geometrically possible sites for an atom of a certain element can be found quickly with minimal compute.
For each of the cliques a node is added to the graph, and then connected to the nodes that make up the clique.

All surface sites can then be classified into unique sites.
This is done by building a sub-graph with a certain number of hops from the clique.
These sub-graphs are then compared to get a list of isomorphs.
The number of hops should be set based on the locality of the chemistry of the surface.
More hops considers longer ranged interactions.
Using the generic FCC(111) surface as an example, there 4 unique surface sites: top, bridge, FCC hollow, and HCP hollow.
For this surface, if 0 hops are selected, the two hollow sites have the same subgraph, and therefore be classified as the same.
If 1 hop is selected, all 4 unique sites are then found; enough neighborhood information is available to distinguish the two sites.
The number of hops is more important for when higher index or irregular surfaces are considered.
Finding these isomorphs is important to reduce workload.
With the classification of isomorphs, only one calculation is required for each isomorph, or more tractably, for each unique site.
These unique sites are classed as anchors.
They are named "anchors" as they are the fixed positions for finding adsorbate sites.
Upon initialization, all reactants are checked for their elements.
For this list of elements, anchors are found across the surface.
As the system is closed, this only has to be computed once.

For input reactants, they are first structurally optimized with a potential.
Again, nodes and edges represent atoms and bonds as before.
In the code implementation, reactants are passed as smiles for ease of use.
Consider the case of the simplest adsorbate: a single atom.
This atom has all possible coordinations already defined from the anchors.
From this, the site stability can be checked with a structure optimization to see if it is a minimum on the potential energy surface (as the KMC should go between minima states).
The final position is then checked for connectivity changes, i.e. bonds formed or broken during the optimization.
If the number of edges is not the same pre and post optimization, the site is unstable.
This allows sites to be pruned for an adsorbate, moving from geometrically possible surface coordinations to physically possible sites.
This allows the quantification of adsorption sites on the surface.

For multi-atom adsorbates, multiple anchors are considered together.
First, an array of all atom-atom distances in the adsorbate is made.
The possible coordinations an adorbate can have to the surface is based on anchor orbits.
To consider the flexibility of a molecule, a nominal buffer is applied to these distances to meet the criteria.
It considers, from an anchor of a designated atom, which other anchors match the distance criteria.
This gives an exhaustive list of geometrically possible adsorbate coordinations to the site.
In addition, it is simplified to unique graph isomorphs  to reduce overheads.
With CO as an example in figXXXX, all possible site coordinations are considered (even those where the CO molecule coordinates to the surface through an oxygen atom!).
This removes any restrictions on the connectivity of the system.
Again, just like the single atom example, structural optimization is performed to find physically possible sites.
In practice, the vast majority of possible coordinations to the surface are pruned.

With surface species now possible, surface reactions can then be considered.
Reactions are considered to be graph transformations.
They can be given strict criteria for specific reactions.
The simplest example is the adsorption or, the reverse process, desorption of an adsorbate.
Physically, it is a site going from an unoccupied to occupied state.
For our implementation, it is the same.
The graph transforms by changing the attribute of an adsorbate site node from unoccupied to occupied, or vice versa.
This reaction is generated when the adsorption site is found.
The occupied and unoccupied energy is calculated when the adsorption site is pruned with a potential and upon the initial catalyst structure optimization.
Free energy corrections can be applied from the harmonic approximation.

The next example of a reaction class is surface diffusion.
This is built from first finding two adsorbate sites of the same adsorbate that share an adsorbate node.
Practically, this is finding two neighboring adsorbate sites.
With this, an initial and final state can be built from the swapping of the site occupation.
As diffusion is a process with a barrier, a transition state must be found between the initial state and the final state.
This is done in the package with the nudged elastic band (NEB) method.

While these "reaction" classes can explain the dynamics of molecules interacting with surfaces, we still do not have reactions taking place on the surface.
An elementary reaction is assumed to be a bond breaking of a surface species, or the reverse, bond forming between two surface species.
This reaction class has the criteria that, for bond breaking, an adsorbate site of a surface species is occupied, and connected to that site are two unoccupied adsorbate sites of the formed fragments when a bond is broken.
For bond formation, the reverse criteria is true.

For bond formation and breaking, which reactions can happen must be discerned.
Bonds are implemented as edges in the graph representation.
With this, a series of bonds can be broken, and this is possible through graph edge deletion.
This forms two components within the graph which can be designated as two new surface species.
This can be done for all bonds in the reactant.
The opposite, bond formation, is much harder to implement.
The package implements it by considering the valency of reactants and surface species.
The graph gives the number of edges of every node, which can infer the valency of all the atoms in the species.
Unsaturated atoms are nodes that can build edges with other unsaturated atoms.
For all species, the atoms that can form bonds are found.
Then, all possible edges that can be formed between two species (where the self interaction is also included) are found.
This is done on the fly.
When a reaction happens and a new surface species is formed, all possible reactions that can happen with that surface species are then calculated.

%Discuss lateral interactions

%Discuss the KMC part

%Do examples

%

%\section{Conclusions}

\begin{acknowledgement}
\end{acknowledgement}

\begin{suppinfo}
\end{suppinfo}

\bibliography{refs}

\end{document}
